% ==========================================================================
%  中英文摘要
% ==========================================================================
\chapter*{\texorpdfstring{摘\quad 要}{摘要}}
\addcontentsline{toc}{chapter}{摘要}
\markboth{摘要}{}

\begingroup
\linespread{1.28}\selectfont

常微分方程初值问题的数值求解是科学与工程计算的基础。当右端函数含有多个时间尺度时，问题呈现刚性：显式方法受稳定性约束，隐式方法每步需要求解非线性方程组。IMEX（隐式—显式）方法对刚性部分采用隐式、对非刚性部分采用显式，在保持稳定性的同时显著降低计算量，是刚性初值问题的一类重要算法。本文以一维粘性 Burgers 方程为例，通过一阶有限体积方法与直线法将其转化为半离散刚性初值问题，系统研究参数化 IMEX-Euler 分裂族的耦合刚性稳定性。

本文引入分裂参数 $\theta\in[0,1]$，将耗散算子按 $\theta$ 与 $1-\theta$ 的比例分配至显式与隐式部分，得到覆盖全显式（$\theta=1$）、对半分裂（$\theta=1/2$）与全隐式（$\theta=0$）的分裂族。针对线性化模型，本文定义双变量稳定性函数
\[
  R(z_E,z_I)=\frac{1+z_E}{1-z_I},
\]
把稳定性问题化为复平面上的一族圆盘与一条耦合刚性射线的位置关系。本文证明：在耗散主导的耦合刚性极限下，$|R|\le 1$ 当且仅当 $\theta\le 1/2$；更精确地，当 $\theta>1/2$ 时，仅当 $|z_D|\le 2/(2\theta-1)$ 时格式稳定。进一步给出连续符号下的联合稳定条件 $\mu^2\le 2d+(1-2\theta)d^2$，并用刚性比 $\varepsilon=z_I/z_E$ 把判据写成 $|\varepsilon|\ge1$。分析表明，只有 $\theta=0$ 是 L-稳定的，$\theta=1/2$ 为临界稳定但不具有刚性衰减。离散 Fourier 符号分析给出了与理论一致的 CFL 型条件，并解释了数值实验中出现的网格尺度不稳定机制。

数值方面，本文在 Burgers 激波算例上完成了四组实验：固定 CFL 扫描黏性系数的刚性相图、不同 $\theta$ 下的最大稳定步长、激波位置与过渡层厚度的定量比较，以及网格收敛性测试。实验结果与理论预测一致：$\theta\le 1/2$ 时格式在刚性极限下保持稳定；$\theta>1/2$ 时最大稳定步长按 $O(\Delta x^2/\nu)$ 收缩，实测与线性预测相差不超过 $2\%$。MUSCL/minmod 重构使 $L^1$ 误差降低 $15$ 至 $28$ 倍，观测收敛阶约为 $2$。

作为扩展，本文将上述语言用于 Burgers 方程对应的 Jin--Xin 松弛系统，比较源项全隐式（策略 A）与仅隐式化线性刚性部分（策略 B）两种分裂。谱分析表明，策略 A 在 $\CFL\le 1$ 时稳定，策略 B 的容许 CFL 上限为 $1/3$，且该限制不随 $\varepsilon\to 0$ 而消失；数值实验验证了策略 B 在 $\CFL=0.7$ 时的失稳。本文的分析框架与代码流程可迁移至更一般的双曲守恒律与刚性松弛系统。

\vspace{1.2em}
\noindent{\bfseries 关键词：}常微分方程初值问题；IMEX 格式；Burgers 方程；耦合刚性；双变量稳定性函数；Jin--Xin 松弛
\endgroup

\clearpage

\chapter*{\texorpdfstring{Abstract}{Abstract}}
\addcontentsline{toc}{chapter}{Abstract}
\markboth{Abstract}{}

\begingroup
\linespread{1.08}\selectfont

Initial value problems (IVPs) for ordinary differential equations (ODEs) are a central object of numerical analysis. When several time scales are present, the IVP is stiff: explicit methods are restricted by stability, while fully implicit methods require a nonlinear solve at every step. Implicit--explicit (IMEX) schemes treat the stiff part implicitly and the non-stiff part explicitly, retaining stability at a reduced computational cost. Taking the one-dimensional viscous Burgers equation as a model, this thesis converts it, by a first-order finite-volume method and the method of lines, into a semi-discrete stiff IVP, and systematically studies the coupled-stiff stability of a parameterised IMEX--Euler family.

A splitting parameter $\theta\in[0,1]$ distributes the diffusion operator between the explicit and implicit parts, covering fully explicit ($\theta=1$), half-and-half ($\theta=1/2$) and fully implicit ($\theta=0$) splittings. For the linearised model, a two-variable stability function
\[
  R(z_E,z_I)=\frac{1+z_E}{1-z_I}
\]
is introduced, which reduces stability to the position of a coupled stiff ray relative to a family of disks in the complex plane. It is proved that $|R|\le 1$ in the diffusion-dominated coupled stiff limit if and only if $\theta\le 1/2$; more precisely, for $\theta>1/2$ stability holds only when $|z_D|\le 2/(2\theta-1)$. A joint stability condition $\mu^2\le 2d+(1-2\theta)d^2$ is derived for the continuous symbol, and the criterion is recast as $|\varepsilon|\ge1$ with the stiffness ratio $\varepsilon=z_I/z_E$. Only $\theta=0$ is L-stable, whereas $\theta=1/2$ is marginally stable but lacks stiff decay. A discrete Fourier symbol analysis yields consistent CFL-type conditions and explains the grid-scale instability observed in the numerical experiments.

Numerical experiments on a Burgers shock include a stiffness phase diagram, maximum stable time steps, and quantitative comparisons of the shock position, the transition-layer thickness and the $L^1$ error. The results agree with the theory: $\theta\le 1/2$ remains stable in the stiff limit, while for $\theta>1/2$ the maximal step scales like $O(\Delta x^2/\nu)$. The MUSCL/minmod reconstruction reduces the $L^1$ error by a factor of $15$--$28$, with an observed order close to two.

As an extension, the same framework is applied to the Jin--Xin relaxation system for Burgers' equation. Two source splittings are compared: a fully implicit one (strategy A) and a linearly implicit one (strategy B). Spectral analysis shows that A is stable for $\CFL\le 1$, whereas B admits only $\CFL\le 1/3$, a restriction that does not disappear as $\varepsilon\to0$; numerical experiments confirm the instability of B at $\CFL=0.7$. The analysis and the code flow are transferable to more general hyperbolic conservation laws with stiff relaxation.

\vspace{1.2em}
\noindent{\bfseries Keywords:} initial value problem for ODEs; IMEX scheme; Burgers equation; coupled stiffness; two-variable stability function; Jin--Xin relaxation
\endgroup
